\section{Threats to validity}
\label{sec:threats}

Each item below names a way the results could be wrong and how much of the paper it affects. The last lists the work that is not finished.

\textbf{Extraction error.} We produced the corpus coding with an LLM-assisted
extraction pass over \NCorpus{} papers and \NCodedFields{} coded fields. Every count in
Sections~\ref{sec:taxonomy} to~\ref{sec:coverage} depends on it. A human validation study is drawn and
stratified by the coding's own confidence: \NValidationDraw{} papers,
\NValidationFields{} fields each, designed for a \NValidationCI{}\pp{}
confidence interval on overall accuracy. It has not been coded yet.
We therefore report no measured extraction-error rate, and the counts in those
sections should be read as unvalidated. The study is single-coded by decision
of the project lead. When it returns it will bound extraction error without
separating a coder's mistake from the extraction's, so it will carry no
inter-rater agreement and we will claim none.

\textbf{Assigning papers to security or safety.} The cross-citation rates depend on classifying every
cited paper as a security or a safety paper. That classification is automated over the full
discovery pool. A systematic bias in it would move the rates. We regard the
direction of the finding as safe, because the gap between \PctEdgeSecCross{}\% of
security papers' citations going to safety papers and \PctSameRegime{}\% of
tested pairs staying inside one literature is far larger than any plausible
classification error. The exact rates should still be read as estimates. Two further limits
bound them. The denominator is the \NAnalysed{} papers that named a
technique rather than every screened-in paper, because a survey cites
widely and names nothing, so counting it would measure how widely surveys cite, not whether methods
cross. And
only citations to papers inside that set register, so a paper reaching
across to work we did not include does not count: every cross-citation rate
in this paper is a floor.

\textbf{Registry granularity.} What counts as one named cause is a
judgement. ``Indirect prompt injection'' is one registry entry and absorbs
\PctTopMechShare{}\% of all tested pairs. \NInjUncovered{} of the \NInjVariants{} named injection
variants have no tested pair at all. A
coarser registry would make coverage look better and a finer one would make it
look worse. We chose the granularity the papers themselves use, which is the
name each paper gives its own contribution.

\textbf{Who-reports-what is observational.} Section~\ref{sec:whoran} reports a
strong association between authorship and outcome. It is not a controlled
comparison. The two sides largely write about different pairs. The
association therefore reflects which pair each side chose to run, rather than
differing conduct on the same pair. That selection is our interpretation.

\textbf{Our experiments are a sample, not a census.} The pairs we ran were
selected by the taxonomy rather than at random, biased toward pairs where
both artifacts run. Fixes whose artifacts do not run are
absent, and nothing here says whether they transfer.

\textbf{Attacks are non-adaptive, with one exception.} This is the criticism
that invalidated the defense rankings of the evaluation platform
DEEPSEC~\cite{carlini-2019-critique}, and
we inherit it. An attack that was not built for a fix and defeats it
anyway bounds attack strength from below, so our failure results are
conservative. A resistance result drawn the same way is not evidence of
robustness at all.
Section~\ref{sec:adaptive} tests the one resistance result we report. Every
other fix we ran would need the same treatment before any claim
that it holds.

\textbf{Variance.} Carlini lists ``no variance
reported'' as one of five fatal flaws in the evaluation platform DEEPSEC~\cite{carlini-2019-critique}.
His own check found a standard deviation of 6.8 points on a reported attack
success rate. That is enough to make many of DEEPSEC's pairwise comparisons
meaningless. We avoid that a
different way: we use greedy decoding everywhere except one cell, so there is
no run-to-run noise to report. The exception is
CheckRLM~\cite{xu-2026-checkrlm}, whose own recipe samples at a temperature
its authors set per backbone. Our standard is to keep an author's
configuration where it differs from ours, so that cell decodes as they
specify and the replication below does not cover it.

Rerunning one published cell that decodes greedily, three
times from scratch with an identical command, produced \textbf{byte-identical
results}: 72 numeric cells across 8 result rows, zero differences, the same
hash after removing the timestamp. Error bars over a deterministic decoder
would be zero-width and would say less than this does. The uncertainty that
remains is sampling uncertainty, which the Wilson intervals~\cite{wilson-1927-probable} and the
detectability threshold already carry.

What this does \emph{not} show is insensitivity to item order, batch
composition, or the few-shot examples a fix holds out. Those are untested,
and a method whose result depended on them would pass this check and still be
fragile.

Nor does it show reproducibility across environments, and we have a
measured example of the difference. The fixes in the retrieval experiments cannot
share one Python environment: the released decoder of the safety fix ACD imports symbols that
newer \texttt{transformers} releases removed, and its authors pin a version
the rest of the cell cannot use. The \emph{clean} arm plants nothing, so its result is a
control. Re-running an earlier ACD cell under a rebuilt stack moved that arm
by up to 0.3\pp{}. The closed-book arm moved too, and
that arm never sees the context at all, so nothing but the environment can
explain it. We therefore re-ran the earlier adaptive passage under the same
stack as the gradient passage rather than comparing across the two; the
attacked arms agreed to within 0.3\pp{} and the corrupted row agreed exactly.
Determinism within a stack and reproducibility across stacks are different
properties, and only the first is established by the replicates above.

\textbf{We do not measure what the fixes cost.} Every number here is an
accuracy, a copy rate or a refusal rate. We report no inference latency,
no memory footprint and no throughput, so nothing in this paper says whether a
fix is affordable to deploy. That matters for the contrastive-decoding
family in particular, which runs two forward passes per token. A method that looks useless here could still be the best choice once
cost is counted, and one that holds up well could be too expensive to
deploy.

\textbf{One dataset and one model per experiment.} The retrieval experiments
run on CoConflictQA~\cite{huang-2025-parammute}, one passage per question,
which is the format of the safety fixes but a different release of Natural
Questions than most of their authors used. Each fix runs on one base model of
7 to 14 billion parameters, and the results should be read as holding for
models of this class and size.

\textbf{Pairs rest on reading.} Every pair counts only after a reader
confirms it against the checklist of Section~\ref{sec:coverage}, and the
readers are LLM-assisted passes over the full text. The checklist is strict:
a result pooled across attacks, or a failure only cited as motivation, does
not count, so the counts are floors. The checklist and every quote ship with the artifact, so a
reader can re-check any pair.

\textbf{The dataset is a pilot.} Its construction rules are not frozen. The
judge's precision falls from \PctSdJudgePrec{}\% on the labels it was tuned on
to \PctSdTwoJudgePrec{}\% on new questions, so the full build needs a
hand-checked sample from every cause. Requiring a clean passage that states
the gold answer drops most questions, which may favor questions with clear
evidence.

\textbf{Scorer.} We reuse each fix's own evaluation rubric where one
exists, and a shared recall-based rubric across the retrieval experiments. Recall
against a gold string rewards a model that produces the right tokens for the
wrong reason. We report the rate of answers from memory and the refusal
rate alongside it so that the failure modes stay visible.

\textbf{One screening stage is not reconstructible.} Figure~\ref{fig:prisma}
counts each stage of corpus construction from recorded provenance. The step
from the \NSurvived{} screened survivors to the \NFullText{} papers read in
full is the exception. The project kept no record of how that subset was
chosen, so we cannot show it was unbiased. A reader should treat the corpus as
a large and systematically discovered sample of the two literatures, not as
their exhaustive intersection. The counts that follow from it are floors.

\textbf{Most of the corpus is unrefereed, and the finding does not depend on
that.} \NVenuePapers{} of \NCorpus{} papers carry a named publication venue;
the rest are preprints, and the share falls with recency, from
\PctVenueEarly{}\% of our \YearVenueEarly{} papers to \PctVenueLate{}\% of
our \YearVenueLate{} ones. A corpus restricted to refereed work
would therefore describe an earlier field than the one we set out to measure,
so we use the refereed subset as a check, not as the corpus. Resolving both sides of
every pair back to a paper succeeds for \NVenueResolved{} of
\NTimelinePairs{} pairs (\PctVenueResolved{}\%); of those, \NVenuePairs{} have
a named venue on both sides, and none of them is a security fix evaluated
against a cause from a safety paper.

\textbf{Census recall.} \NNoTechnique{} papers contributed no cause and no fix, and
\NNoTechOther{} of those match none of the survey, benchmark, framework or
position descriptions by title. We re-read \NRecallRead{} papers --- that
bucket as it stood before the pass --- and asked one question of each: does
this paper name a technique we failed to register. \NRecallCandidates{}
candidates came back. Two were genuine omissions and are now in the census. One was
a fix already registered under a different name. One renamed a cause
the registry carries three times. One was a form of planted passage rather than a
technique. The last was an attack a defense paper built for its own
evaluation; Section~\ref{sec:literatures} registers it under the rule that
bars a paper's own attack from pairing with its own defense. Two omissions in
\NRecallRead{} papers is a rate, not a guarantee of completeness.

\textbf{Absence is not measured.} Every count in Section~\ref{sec:coverage}
is a positive finding, something a paper says, not something a pass failed to
find. We built three automated passes over these papers, and hand-checking
each found the same asymmetry: what a pass flagged as present was reliable,
and what it flagged as absent was about half wrong. A pattern can confirm
that something is present in a paper; it cannot establish that something is
absent, because absence of the phrasing it knows is not absence of the thing.
We therefore report no count of the second kind.

\textbf{To be done.} Three items are not complete, and we report none of them
as complete.

\begin{itemize}
\item \textbf{Code the validation study.} The draw is complete. We give no
  extraction-error rate until the study returns.
\item \textbf{Read two published benefits from their papers.} The safety
  fixes ACD and DyPRAG give no figure in their text, so their controls stay
  directional.
\item \textbf{Run both directions on the shared dataset.}
  Until then the reverse direction, security defenses against honest
  mistakes, is untested here.
\end{itemize}
